\section{Proofs for Expressiveness Lower Bounds (\Cref{sec:lower_bounds})}\label{appsec:proofs}

\newcommand{\interleave}[2]{{#1}^\frown{#2}}
We first introduce some notations. Since in the construction of the lower bounds we are only using fixed-point numbers, we will use the shorthand $\Floating_{s}\triangleq \Floating_{0,s}=\{c\cdot k\cdot 2^{-s}\mid c\in \{-1,1\}, 0\le k\le 2^{2s}-1, k\in\mathbb{N}\}$ and rounding operation $\rds{\cdot} \triangleq[\cdot]_{0,s}$. We use $1_s$ to denote all-one vectors of length $s$. Similarly we define $\inner{\cdot}{\cdot}_s$, $\times_s$, and $\softmax_s$. We recall that for any $s\in \mathbb{N}^+$ and integer $0\le x\le 2^s-1$, we use $\bin_s(x)\in\{0,1\}^s$ to denote the usual binary encoding of integer $x$ using $s$ binary bits in the sense that $x = \sum_{i=1}^s 2^i (\bin_s(x))_i$ and $\sbin_s(x)\in\{-1,1\}^s$ to denote the signed binary encoding, which is $2\bin_s(x)-(1,\ldots,1)$. 



%
%\begin{proof}[Proof of \Cref{lem:cot_1_is_squal_to_O1}]
%	This is a special case of \Cref{thm:cot_dpolyn_slogn_main}.
%\end{proof}

Recall $B_s = \max \Floating_s = 2^s-2^{-s}$.
\begin{lemma}\label{lem:exp_rounding}
	For any $s\in\mathbb{N}^+$, it holds that $\rds{\exp(-B_s)} = 0$.
\end{lemma}
\begin{proof}[Proof of \Cref{lem:exp_rounding}]
	By the definition of rounding operation for $2s$-bit precision (\Cref{defi:rounding}), it suffices to show that $\exp(-B_s)\le 2^{-s-1}$, that is, $B_s \ge \ln 2\cdot (s+1)$. Note that $2^s\ge s+1$ for all $s\ge 1$, we have $B_s/(s+1)\ge B_s2^{-s} =1-2^{-2s}\ge 3/4 \ge \ln 2$.
\end{proof}

Using the same argument above, we also have \Cref{lem:exp_rounding_up}.
\begin{lemma}\label{lem:exp_rounding_up}
	For any $s\in\mathbb{N}^+$, it holds that $\rds{\exp(B_s)} = B_s$.
\end{lemma}

%Finally, we use $\pos_s(x)$ to denote the interleaving between $\sbin_s(x)$ and all-one vector of the equal length, that is, $\pos_s(x)$
\subsection{Proof of \Cref{thm:cot_dlogn_main}}\label{appsec:proof_cot_dlogn_main}
Given two vectors $x,y$ of the same length $s$, we use $\interleave{x}{y}$ to denote their interleaving, that is, $(\interleave{x}{y})_{2i-1} = x_i, (\interleave{x}{y})_{2i} = y_i$ for all $i\in [e]$. 
\begin{lemma}\label{lem:attention_rounding}
	For any $s\in\mathbb{N}^+$, let $q_i = \interleave{\sbin_s(i)}{1_s}$ and $k_i = B_s\cdot (\interleave{\sbin_s(i)}{(-1_s)})$ for all $i\in [2^s-1]$, it holds that $\rds{\exp(\inner{q_i}{k_j}_s)}=\one[i=j]$ for all $i,j\in [2^s-1]$.
\end{lemma}

\begin{proof}[Proof of \Cref{lem:attention_rounding}]
	It suffices to prove that $\inner{q_i}{k_j}_s = -B_s$ if $i\neq j$ and $\inner{q_i}{k_j}_s = 0$ if $i= j$. The rest is done by \Cref{lem:exp_rounding}.
	
	Given any $i,j\in [2^s-1]$, by definition of finite-precision inner product, we know that for any $l\in [2s-1]$,  it holds that $a_{l} = \rds{	a_{l-1} + \rds{(q_i)_{l}(k_j)_{l}}}$ where $a_0 \triangleq 0$ and $a_{l}\triangleq \inner{(q_i)_{:l}}{(k_j)_{:l}}_s$ for $l\in [2s]$. 
	
	For all $l\in [e]$, we have that $\rds{(q_i)_{2l}(k_j)_{2l}} = -B_s$ and $\rds{(q_i)_{2l-1}(k_j)_{2l-1}} = B_s\cdot \ind[(\sbin_s(i))_l = (\sbin_s(j))_l]$. If $i=j$, it is straightforward that $a_{2l-1} = -B_s$ and $a_{2l} = 0$ for all $l\in[e]$. If $i\neq j$, then there exists $l\in[s-1]$ such that $(\sbin(i)))_l\neq (\sbin(j)))_l$. Thus $\rds{(q_i)_{2l-1}(k_j)_{2l-1}} = \rds{(q_i)_{2l}(k_j)_{2l}} = -B_s$, which implies $a_{2l} = -B_s$ regardless of the value of $a_{2l-2}$. Again use induction we can conclude that $a_{2l'} = -B_s$ for all $l\le l' \le e $.
\end{proof}


\begin{proof}[Proof of \Cref{thm:cot_dlogn_main}]
	For any $\gL\in \SIZE[T(n)]$, by definition there is a family of boolean circuit $\{C_n\}$ which compute $\gL$ for all inputs of length $n$ using $O(T(n))$ many   $\NOT$ and $\AND$ gates. Without loss of generality, let us assume the number of non-input gates in $C_n$ be $T(n)$. We will show that for each $n$, there is a 2-layer decoder-only transformer $\transformer_{\theta_n}$ that computes $C_n(x)$ using $T(n)$ steps of CoT for all $x\in\{0,1\}^n$. More precisely, we will construct a transformer that simulates one boolean gate in $C_n$, following the topological order of the circuit, in each step of its chain of thought. 
	
	We first index the gates (including input gates) from $1$ to $n+T(n)$ according to the topological order. For each gate $i \in [n+1,n+T(n)]$, we use $a(i)$ and $b(i)$ to denote its two input gates. Since $\NOT$ only has one input, we set $a(i)$ as its input and $b(i)$ as $0$.  We let $c(i)=0$ if $i$th gate is $\NOT$ and  $c(i)=1$ if $i$th gate is $\AND$. For any input gate $1\le i\le n$,  $a(i),b(i)$, and the gate type are not meaningful and their choice will not affect the output and thus can be set arbitrarily. For convenience, we will set $a(i)=b(i)=c(i)=0$.  We use $g_i(x)$ to denote the output of non-input gate $i$ ($n+1\le i\le n+T(n)$) on the circuit input $x\in\{0,1\}^n$, which is equal to $(1-c(i))(1-x_{a(i)}) + c(i)(\relu(x_{a(i)} + x_{b(i)}-1))$.
	
	Now we describe the construction of the vocabulary $\gV$, the token embedding $\theta_{\tokenembedding}$,  and position encoding $\theta_\posencoding$. We set precision $s$ as any positive integer larger than $1$, $\gV=\{0,1\}$, $k\triangleq k(n) = \lceil \log_2 (T(n)+n) \rceil = O(\log n)$ since $T(n)$ is a polynomial, $d'(n) = 3k+6$, $\theta_{\tokenembedding}(0) = 0\cdot e_1$,  $\theta_{\tokenembedding}(1) =  1\cdot e_1$, and 
	\begin{align}
	(\theta_{\posencoding})_{i-1}^\top = [0,0,0,0, c(i), \sbin_k(i-1) , \sbin_k(a(i)), \sbin_k(b(i)),1], \quad  \forall 2\le i \le n+T(n).	
	\end{align}
	
	We use $h^0_i$, $h^{0.5}_i$, $h^1_i$, $h^{1.5}_i$, and $h^2_i$ to denote the intermediate embeddings at position $i$ and different depths. Here, depth $0.5$ and $1.5$ refer to the output of the Attention layer inside each transformer layer. 
	\begin{enumerate}[leftmargin=*] 
		\item 	For the first attention layer, denoting embedding at the $i$th position $h_i^0$ by $h$, we set the query as $q_i \triangleq \left(h_{k+6:2k+5}\right)^\frown \left(h_{3k+6}\cdot 1_s\right)$, the key as $k_i\triangleq \left(B_sh_{6:k+5}\right)^\frown \left(-h_{3k+6}\cdot 1_s\right)$, and the value as $v_i\triangleq h_1 \cdot e_2$. \footnote{Note here the dimension of $k_i$ and $q_i$ are the same but less than $d'$, which does not strictly satisfy our definition of transformer in \Cref{alg:defi_attn}. This is for notational convenience and is without loss of generality because we can pad extra zeros.}
		\item For the first fully-connected layer, we skip it by setting its weights to be $0$.
		\item For the second attention layer, denoting embedding at the $i$th position $h_i^1$ by $h$, we set the query as $q_i \triangleq \left(h_{2k+6:3k+5}\right)^\frown \left(h_{3k+6}\cdot 1_s\right)$, the key as $k_i\triangleq \left(B_sh_{6:k+5}\right)^\frown \left(-h_{3k+6}\cdot 1_s\right)$, and the value as $v_i\triangleq h_1 \cdot e_3$. 
		\item For the second fully-connected layer, we define $F(a,b,c)\triangleq \relu(1-a-c) + \relu(a+b+c-2)$. Denote  the embedding at position $i$, $h_{1.5}^i$ by $h$. The output of the second fully-connected layer is defined as $F(h_2, h_3, h_4)\cdot e_4$. Note this expression is valid because it can be expressed by two-layer ReLU nets with constant bits of precision and a constant number of neurons.
%		It is easy to see that for $a,b,c\in\{0,1\}$, $F(a,b,c) = 1-a = \NOT a$ if $c=0$ and $F(a,b,c) = a \AND b = \relu(a+b-1)$ if $c=1$. 
		\item The final output at position $i$ is $(h_i^2)_4$. 
	\end{enumerate}
	
	Below we first describe the value of the internal variables of the transformer and then show there exist parameters making such computation realizable. Let $(x_1,\ldots, x_n)$ be the input tokens, we define $x_{n+i}\triangleq \transformer^i_{\theta_n}(x_1,\ldots,x_n), \forall 1\le i\le T(n)$. We claim there exists transformers parameter $\theta_n$ such that $\transformer^{T(n)}_{\theta_n}(x_1,\ldots,x_n) = g_{T(n)}(x)$ (\emph{i.e.} $C_n(x)$). More specifically, we claim that our constructions will ensure the following inductively for all $n+1\le i\le n+T(n)$,
	\begin{enumerate}
		\item $h^0_{i-1}= [x_{i-1},0,0,0, c(i), (\sbin_k(i-1))^\top , (\sbin_k(a(i)))^\top, (\sbin_k(b(i)))^\top,1]^\top$;
		\item $h^{0.5}_{i-1}\!= [x_{i-1},x_{a(i)},0,0, c(i),(\sbin_k(i-1))^\top , (\sbin_k(a(i)))^\top, (\sbin_k(b(i)))^\top,1]^\top$;
		\item $h^{1}_{i-1}= [x_{i-1},x_{a(i)},0,0, c(i),(\sbin_k(i-1))^\top , (\sbin_k(a(i)))^\top, (\sbin_k(b(i)))^\top,1]^\top$;
		\item $h^{1.5}_{i-1}\!\!=\!\! [x_{i-1},x_{a(i)},x_{b(i)},0, c(i),(\sbin_k(i-1))^\top\!, (\sbin_k(a(i)))^\top\!, (\sbin_k(b(i)))^\top,1]^\top$;
		\item $h^{2}_{i-1}\!\!= \!\![x_{i-1},x_{a(i)},x_{b(i)}, g_i(x),c(i),(\sbin_k(i-1))^\top , (\sbin_k(a(i)))^\top, (\sbin_k(b(i)))^\top,1]^\top$;
		\item $x_{i} \triangleq \transformer_{\theta_n}(x_1,\ldots,x_n,\ldots,x_{i-1}) =  g_{i}(x)$.
	\end{enumerate} 
	
	Now we explain why the above conditions hold for any position $i$ using induction, \emph{i.e.}, assuming it is true for all $n+1\le i'\le i$. We first notice that by our construction, for all $1\le i' \le n-1$, it holds that $(h_{i'}^k)_1 = x_{i'}$ and $(h_{i'}^k)_{6:k+5}$ for all $k\in\{0,0.5,1,1.5,2\}$. Note these are the only information that will be used in the later attention layers.
	\begin{enumerate}
		\item This is simply by the construction of $\theta_{\tokenembedding}$ and $\theta_{\posencoding}$.
		\item In the first attention layer, at the $j$th position, we have $q_j \triangleq \sbin_k(a(j))^\frown 1_s$ as the query, $k_j \triangleq B_s \cdot \sbin_k(j)^\frown (-1_s)$ as the key and $v_j\triangleq x_j$ as the value for all $j\in[n+T(n)]$. Note here we reduce the sizes of hidden embeddings for simplicity of demonstration. This is valid because we can fill the extra coordinates by 0. This is valid because we can always set the extra coordinates to be $0$. By \Cref{lem:attention_rounding}, we know that $\rds{\exp(\inner{q_i}{k_j}_s)}=\one[a(i)=j]$ for all $j\in[n + T(n)]$. Recall that the attention scores are defined as $s_i \triangleq \softmax(\inner{q_i}{k_1}_s,\ldots,\inner{q_i}{k_i}_s) \| (0,\ldots, 0) $, we know that $s_i = e_{a(i)}$.
		
		 
		
		\item We set the parameters in the first fully-connected feedforward layer to be all $0$ and let the skip connection pass the intermediate values.
		\item The second attention layer attains $x_{b(i)}$ and places it in the third coordinate in the same way as step 2. 
		\item In the fully-connected feedforward layer we compute $F(x_{a(i)}, x_{b(i)}, c(i)) = \relu(1 - x_{a(i)}-c(i)) + \relu(x_{a(i)} + x_{b(i)} + c(i)-2)$ for all $x_{a(i)},x_{b(i)},c(i)\in\{0,1\}$. We can verify that $F(x_{a(i)}, x_{b(i)}, c(i))= (1-c(i))(1-x_{a(i)}) + c(i)(\relu(x_{a(i)} + x_{b(i)}-1)) = g_i(x)$,  which is the desirrd output of the gate $i$. This is because when $c(i)=0$, the output is $1-a(i)=\NOT a(i)$ and when $c(i)=0$, the output is $\relu(a(i)+b(i)-1) = a(i) \AND b(i)$.  
		\item The output layer uses the fourth coordinate of $h^2_i$, which is $g_i(x)$ according to induction, as the output. 
	\end{enumerate}
This completes the proof of \Cref{thm:cot_dlogn_main}.
\end{proof}


\subsection{Proof of \Cref{thm:cot_dpolyn_slogn_main,thm:cot_dpolyn_s1_main}}\label{appsec:proof_cot_dpolyn_slogn_main}

In this subsection, we prove \Cref{thm:cot_dpolyn_slogn_main,thm:cot_dpolyn_s1_main}. We first prove a useful lemma that gives an equivalent characterization of $\SIZE^{\AC^0}$ and $\SIZE^{\TC^0}$.

\begin{lemma}\label{lem:size_Tn_with_TC0_oracle}
	For any $T(n)\in \poly(n)$ satisfying $T(n)\ge 1,\ \forall n\in\mathbb{N}^+$, a  decision problem $\gL:\cup_{k\in \mathbb{N}^+}\{0,1\}^k\to \{0,1\}$ belongs to $\SIZE^{\AC^0}[T(n)]$ (resp. $\SIZE^{\TC^0}[T(n)]$)  if and only if  there exist a polynomial $S(n)$, a function $T'(n)=O(T(n))$, and a depth $L\in\mathbb{N}^+$  such that for every $n\in\mathbb{N}^+$ there exist a sequence of sizes-$S(n)$, depth-$L$ circuits, $\{C_n^{i}\}_{i=1}^{T'(n)}$, with unlimited-fanin $\AND$, $\OR$ and $\NOT$ gates (with additionally $\MAJORITY$ gates for $\SIZE^{\TC^0}$) and that for all $ x\in\{0,1\}^n$,
	\begin{align}
		\gL(x) = x_{n+T'(n)}, \text{\quad where \quad }  x_{n+i}\triangleq C_n^i(x_1,\ldots, x_{n+i-1}), \quad \forall i \in [T'(n)].
	\end{align}
\end{lemma}

\begin{proof}[Proof of \Cref{lem:size_Tn_with_TC0_oracle}]
We will prove for $\SIZE^{\TC^0}$ only and the proof for $\SIZE^{\AC^0}$ is almost the same. 
	
The ``$\implies$'' direction is straightforward. By definition of $\SIZE^{\TC^0}[T(n)]$~(\Cref{defi:circuits_with_oracle}), for any $\gL\in \SIZE^{\TC^0}[T(n)]$, there is a function $p(n)\in\poly(n)$ and a family of $\TC^0$ circuits $\{C'_i\}_{i=1}^\infty$ such that for every $n\in \mathbb{N}$ and $x\in\{0,1\}^n$, $\gL_n(x_1,\ldots,x_n)$ can be computed by a size-$O(T(n))$ threshold circuits with oracle gate $C_{p(n)}$. Now we sort all the nodes in the circuits with oracle gates along the topological order as $x_1,\ldots,x_{n+T'(n)}$ where $x_1,\ldots,x_n$ are the inputs and $T'(n)=O(T(n))$ is the number of the gates, then clearly $x_{n+i}$ is a function of $x_1,\ldots,x_{n+i-1}$ for all $i\in[T'(n)]$ and this function can be implemented by a different threshold circuit of constant depth and $\poly(n)$ size for each $i$. This completes the proof of ``$\implies$'' direction.
	
Now we prove the other direction ``$\Longleftarrow$''. We first show that given $T'(n)$ sizes-$S(n)$, depth-$L$ circuits, $\{C_n^{i}\}_{i=1}^{T'(n)}$, there is a depth-$(L+1)$, size $O(T'(n)S(n))$ circuit $C'_n$, such that 
\begin{align}
C'_n(x_1,\ldots, x_{n+T'(n)-1},e_{j}) = C_n^j(x_1,\ldots,x_{n+j-1}),\quad \forall j\in [T'(n)], x\in\{0,1\}^{n+T'(n)-1},
\end{align}
where $\one_j\in\{0,1\}^{T'(n)}$ is the one-hot vector with its $j$th coordinate being $1$. Indeed, it suffices to set 
\begin{align}
	C'_n(_1,\ldots, x_{n+T'(n)-1}, y_1,\ldots, y_{T'(n)})  = \vee_{j=1}^{T'(n)} \left( y_j \wedge C_n^j(x_1,\ldots,x_{n+j-1}) \right).
\end{align}


Once we have such oracle gate $C'_n$, given input $x_1,\ldots,x_n$, we can recursively define 
\begin{align}
x_{n+i} \triangleq C'_n(x_1,\ldots, x_{n+i-1}, 0_{T'(n)-n}, e_{j}).	
\end{align}
Thus we can compute $\gL(x)=x_{n+T'(n)}$ using $T'(n)$ oracle gate $C'_n$. We can get constant gate $0$ and $1$ by using $x_1 \wedge \neg x_1$ and $x_1 \vee \neg x_1$. respectively. This completes the proof. 
\end{proof}

Now we are ready to prove \Cref{thm:cot_dpolyn_slogn_main,thm:cot_dpolyn_s1_main}. We will prove \Cref{thm:cot_dpolyn_slogn_main} first and the proof of \Cref{thm:cot_dpolyn_s1_main} is very similar to \Cref{thm:cot_dpolyn_slogn_main}.


\begin{proof}[Proof of \Cref{thm:cot_dpolyn_slogn_main}]

We first show that $\SIZE^{\TC^0}[T(n)+1] \supseteq \Cot[T(n)][ \poly(n)][\log n]$. For the case that the vocabulary of transformer $\gV=\{0,1\}$, by \Cref{thm:TC0_upper_bound}, we know for any $\theta_n$, $\transformer_{\theta_n}(x_1,\ldots,x_i)$ can be expressed by a $\TC^0$ circuit whose depth is uniformly upper bounded by some constant for all $n\le i\le n+O(T(n))$. This completes the proof when $\gV=\{0,1\}$. When $\gV\neq \{0,1\}$, we can use the binary encoding of elements in $\gV$ as the inputs of those $\TC^0$ gates constructed for the later layers of the transformer. 

Now we turn to the proof for the other direction: $\SIZE^{\TC^0}[T(n)+1] \subseteq \Cot[T(n)][ \poly(n)][\log n]$. In high-level speaking, the proof contains two steps: 
\begin{enumerate}
	\item  We show that $\TC^0 \subseteq \T[\poly(n)][\log n]\subseteq \Cot[1][ \poly(n)][\log n]$. The first step has two key constructions: (a). using attention to copy all the weights to the same position; (b). we can use polysize two-layer FC net with ReLU activation to simulate $\MAJORITY,\AND,\OR$ gate with unbounded fan-in~(\Cref{lem:FB_simulating_boolean_gates});
	\item We can do the first step for all positions $i=n+1,\ldots,n+O(T(n)+1)$ simultaneously.
\end{enumerate}



By the equivalence construction shown in the \Cref{lem:size_Tn_with_TC0_oracle}, we know that for any problem $\gL\in \SIZE^{\TC^0}[T(n)+1]$, there exist constant $L$, polynomial $S(n)$, and $T'(n)=O(T(n)+1)$, and a sequence of threshold circuits, $\{C_n\}_{n\in\mathbb{N}}$, each of size $S(n)$ (number of non-input gates) and depth of $L$, and that for all $ x\in\{0,1\}^n$,
	\begin{align}
		L(x) = x_{n+T'(n)}, \text{\quad where \quad }  x_{n+i}\triangleq C_n(x_1,\ldots, x_{n+i-1},\underbrace{0,\ldots,0}_{T'(n)-i},e_i), \quad \forall i \in [T'(n)].
	\end{align}
	
	Now we present the construction of the constant-depth, constant-precision decoder-only transformer, $\transformer_{\theta_n}$ which computes problem $\gL$ when input length is $n$. Without loss of generality, we only consider the case where $T'(n) = T(n)+1$. We set vocabulary $\gV = \{0,1\}$, embedding width $d(n) = 1 + 3(T(n)+n)+S(n) = O(\poly(n))$, depth equal to $L+1$, CoT length $T(n)$ and precision $s(n)= \lceil \log_2 S(n)\rceil$ so the precision is high enough for simulating all the $\poly(n)$ size $\MAJORITY$ gates used in $C_n$ (\Cref{lem:FB_simulating_boolean_gates}). We set $(\theta_{\tokenembedding})_0 = 0\cdot e_1$,  $(\theta_{\tokenembedding})_1 =  1\cdot e_1$, and $(\theta_{\posencoding})_i = e_{i+1}$ for all $i\in[n+T(n)]$, where we use $e_i\in \{0,1\}^{d(n)}$ to denote the one-hot vector whose $i$th coordinate is $1$ for $i\in[d(n)]$ and $\overline e_i\in\{0,1\}^{n+T(n)}$ to denote one-hot vector whose $i$th coordinate is $1$ for $i\in[n+T(n)]$. 
	
Below we first describe the value the internal variables of the transformer and then show there exist parameters making such computation realizable. To make our claims more interpretable, we only write the non-zero part of the embedding and omit the remaining $0$'s.  the Let $(x_1,\ldots, x_n)$ be the input tokens and $\Delta\triangleq 4n+4T(n)+1$, our constructions will ensure that 
	\begin{enumerate}
		\item $x_{n+i}= \transformer^i_{\theta_n}(x_1,\ldots,x_n)$,  $\forall i\in[T(n)]$. 
		\item $h^0_{i}= x_i e_1 + e_{i+1} = (x_i, \overline e_i) \  \forall i \in [ n+T(n)]$;
		\item $h^{0.5}_{i}\!= x_i e_1 + e_{i+1} = (x_i, \overline e_i), \  \forall i \in [ n+T(n)]$;
		\item $h^{1}_{i} = x_i e_1 + e_{i+1} + x_ie_{n+T(n)+i+1} = (x_i, \overline e_i, x_i \overline e_i), \  \forall i \in [ n+T(n)]$;
		\item $h^{1.5}_{i} =  (x_i, \overline e_i, x_i \overline e_i, x_1,x_2,\ldots,x_i), \forall i \in [ n+ T(n)]$
		\item $h^{2}_{i} =  (x_i, \overline e_i, x_i \overline e_i, x_1,1,x_2,1,\ldots,x_i,1), \forall i \in [ n+ T(n)]$
		\item $(h^{1+l}_{n+i-1})_{\Delta+1: \Delta+ 2S(n) }\in\{0,1\}^{S(n)}$ stores the intermediate result of circuit $C_n^{i}(x_1,\ldots,x_{n+i-1}, 0,\ldots, 0,,e_i)$ at layer $l$, $\forall i\in[T(n)]$ and $l\in [L]$;
		\item $(\theta_{\transoutput} h^{L+2}_{n+i-1})_0 = 1-2C_n^{i}(x_1,\ldots,x_{n+i-1}, 0,\ldots, 0,e_i), (\theta_{\transoutput} h^{L+2}_{n+i-1})_1 =0$, for all $i\in[T(n)]$.
	\end{enumerate} 
	
	
	Now we explain the purpose of each layer and how to set the parameters such that the requirements above are met.
	\begin{enumerate}
		\item $x_{n+i}= \transformer^i_{\theta_n}(x_1,\ldots,x_n)$,  $\forall i\in[T(n)]$ is the goal of the construction;
		\item This is by our construction of $\theta_\posencoding$ and $\theta_\tokenembedding$;
		\item The first attention layer does nothing by setting all weights to $0$;
		\item By \Cref{lem:FB_simulating_boolean_gates}, $\AND$ can be simulated by 2-layer ReLU networks using $2$ hidden neurons.  Thus we use the first feedforward-layer to compute the function $(h_i^1)_{n+T(n)+1+j} =  (h_i^{0.5})_{1+j}\  \wedge \ (h_i^{0.5})_{1}$ for all $i,j\in[n+T(n)]$ with totally $2(n+T(n))$ hidden neurons. Therefore if $j\neq i$, then $ (h_i^{0.5})_{1+j} = 0$, which implies $(h_i^1)_{n+T(n)+1+j}  = 0$; if $j=i$, then $(h_i^{0.5})_{1+j} = 1$, thus $(h_i^1)_{n+T(n)+1+j}(h_i^{0.5})_{1} = x_i$. 

		\item This step exactly requires $(\Attn_{\theta^{(1)}_{\Attn}}(h^{(1)}_1,\ldots,h^{(1)}_n))_{i} = \sum_{j=1}^i e_{\Delta+j}x_j$. It suffices to set the attention score of the second layer at $i$th position  $s_i= (1,\ldots,1,0\ldots,0) = (1_i,0_{n+T(n)-i})$ for all $i\in[n+T(n)-1]$. This can be done by setting $q^{1.5}_i = W_Q h^{1.5}_i = (B_{s(n)}, 0_{n+T(n)-1}), k^{1.5}_i = W_K h^{1.5}_i = (1,0_{n+T(n)-1})$. By \Cref{lem:exp_rounding_up}, we have $[\exp([\inner{q_i}{k_j}]_{s(n)})]_{s(n)}=[\exp(B_{s(n)})]_{s(n)} = B_{s(n)}$. Since rounded sum of any number of $B_{s(n)}$ is still $B_{s(n)}$ and $[B_{s(n)}/B_{s(n)}]_{s(n)}=1$, we know that  
		$$s_i= \softmax_{s(n)}(B_{s(n)}1_{i}) \| (0,\ldots,0) = (1,\ldots,1,0\ldots,0) = (1_i,0_{n+T(n)-i})$$ 
		for all $i\in[n+T(n)-1]$. Note in this step we use our specific rounding rule to copy all the previous $x_i$ with a sum of attention score larger than $1$. We can just also use approximately uniform attention scores with an additional coefficient before $x_i$ since we have $\log n$ precision.
		Finally we set $v^{1.5}_i = W_V h^{1.5}_i = e_{\Delta+i}x_i$ and $W_O = I_{d(n)}$.
	\item The second MLP layer just does permutation and adds some constants into fixed coordinates. The construction is straightforward and thus omitted.
	\item The second attention layer is the only attention layer which has non-zero weights. Using the feedforward ReLU networks from layer $2$ to $L+1$, we can simulate the circuits $C_n$ in parallel for all $i\in[T(n)]$ by \Cref{lem:FB_simulating_boolean_gates}. In detail, \Cref{lem:FB_simulating_boolean_gates} ensures that we can use a two-layer fully-connected ReLU network with weights  to simulate a layer of the $\TC^0$ circuits $C_n$. Moreover, there is enough space in the embedding to reserve $S(n)$'s $1$ needed by \Cref{lem:FB_simulating_boolean_gates}. 

	\item This step holds directly due to the property guaranteed in step $8$. We note that with the property claimed in step 9, we have that $(\theta_{\transoutput} h^{L+2}_{n+i-1})_1 - (\theta_{\transoutput} h^{L+2}_{n+i-1})_0 = 2C_n(x_1,\ldots,x_{n-i+1}))-1$. Thus if $C_n(x_1,\ldots,x_{n-i+1}) =1$, then $(\theta_{\transoutput} h^{L+2}_{n+i-1})_1 - (\theta_{\transoutput} h^{L+2}_{n+i-1})_0 = 1$, which implies $\transformer_{\theta_n}(x_1,\ldots,x_{n+i-1}) = 1$, otherwise if $C_n(x_1,\ldots,x_{n-i+1}) =0$, then $\transformer_{\theta_n}(x_1,\ldots,x_{n+i-1}) = 0$. In both cases, we have that 
		\begin{align}
			C_n(x_1,\ldots,x_{n-i+1}) = \transformer_{\theta_n}(x_1,\ldots,x_{n+i-1})
		\end{align}
	\end{enumerate}
	So far we have finished the proof for the general $T(n)$. Specifically, when $T(n)=T'(n)=0$,  our proof shows that the constant-depth transformer can still simulate any constant-depth circuit, which means $\TC^0\subseteq \T[\poly(n)]\subseteq \Cot[1][\poly(n)] = \SIZE^{TC^0}(1) = \TC^0$. Thus all the inclusions are equivalence, that is $\TC^0= \T[\poly(n)]= \Cot[1][\poly(n)] = \SIZE^{TC^0}(1)$.
\end{proof}


\begin{proof}[Proof of \Cref{thm:cot_dpolyn_s1_main}]

We first show that $\SIZE^{\AC^0}[T(n)+1] \supseteq \Cot[T(n)][ \poly(n)][1]$. For the case that the vocabulary of transformer $\gV=\{0,1\}$, by \Cref{thm:AC0_upper_bound}, we know for any $\theta_n$, $\transformer_{\theta_n}(x_1,\ldots,x_i)$ can be expressed by a $\TC^0$ circuit whose depth is uniformly upper bounded by some constant for all $n\le i\le n+O(T(n))$. This completes the proof when $\gV=\{0,1\}$. When $\gV\neq \{0,1\}$, we can use the binary encoding of elements in $\gV$ as the inputs of those $\TC^0$ gates constructed for the later layers of the transformer. 

The other direction is almost the same as that of \Cref{thm:cot_dpolyn_slogn_main}, except that we now only need constant bits of precision because we do not need to simulate $\MAJORITY$ gates (\Cref{lem:FB_simulating_boolean_gates}).

\end{proof}

\subsection{Proof of \Cref{thm:transformer_polyn_stronger_than_logn}}\label{appsec:proof_transformer_polyn_stronger_than_logn}
\begin{proof}[Proof of \Cref{thm:transformer_polyn_stronger_than_logn}]
By \Cref{lem:TC0_not_in_fixed_poly}, it holds that for all $k\in\mathbb{N}$, $\AC^0\subsetneq \SIZE[n^k]$. By \Cref{thm:cot_dpolyn_slogn_main}, we know that $\AC^0\subseteq \Cot[1][\poly(n)][1] \subseteq \Cot[n^k][\poly(n)][1]$ for any $k\in\mathbb{N}$. Thus $\Cot[n^k][\poly(n)][1]\subsetneq \SIZE[n^k]$ for all $k,k'\in\mathbb{N}$. Also, note that the attention layer and fully-connected layer can be computed using poly-size circuits. Thus for any $k\in\mathbb{N}$, $\Cot[n^k][\log(n)]\subseteq \SIZE[n^{k'}]$ for some integer $k'\ge k$. Combining these we conclude that for any $k\in\mathbb{N}$, $ \Cot[n^k][\log(n)]\subsetneq \Cot[n^k][\poly(n)]$.
\end{proof}

\subsection{Auxiliary Lemmas}\label{appsec:auxiliary_lemmas}

In this subsection, we prove a few auxiliary lemmas that are used in the proofs in \Cref{sec:lower_bounds}.

\begin{lemma}\label{lem:FB_simulating_boolean_gates}
Unlimited-fanin $\AND,\OR$ (resp. $\MAJORITY) :\{0,1\}^n\to \{0,1\}$ can be simulated by some 2-layer feedforward ReLU network with constant (resp. $\log n$) bits of precision constant hidden dimension and additional $n$ constant inputs of value 1. 

Mathematically, let $\FF[s(n)]$ be the set of functions $C:\{0,1\}^n\to\{0,1\}$ which can be a two-layer feedforward ReLU network with at most $s(n)$ bits of precision and constant hidden dimension $\FF_{\theta}:\{0,1\}^{2n}\to \{0,1\}, \FF_{\theta}(x') = W_2\times_s\relu(\rds{W_1\times_s x'+b_1})$, where $\theta = (W_2,W_1,b_1)$, such that for any $x\in\{0,1\}^n$, 
\begin{align}
\FF_\theta(x_1,1,x_2,1,\ldots,x_n,1) = C(x).	
\end{align}
We have unlimited-fanin $\AND,\OR\in \FF[1]$ and $\MAJORITY\in \FF[\log n]$.  
\end{lemma}

The proof of \Cref{lem:FB_simulating_boolean_gates} is based on the following straightforward lemma~(\Cref{lem:indicator}).
\begin{lemma}\label{lem:indicator}
	For any $s\in\mathbb{N}^+$ and $a\in\mathbb{Z}\cap \Floating_s$, $\relu(\rds{a})-\relu(\rds{a-1}) = \ind[a>0]$. 
	In particular, for any $a\in\mathbb{Z}$, $\relu(a)-\relu(a-1) = \ind[a>0]$.
\end{lemma}


\begin{proof}[Proof of \Cref{lem:FB_simulating_boolean_gates}]


	Recall that $\interleave{x}{y}$ denotes $(x_1,y_1,x_2,y_2,\ldots,x_n,y_n)$ for any $x,y\in\{0,1\}^n$. We have that $\sums(\interleave{x}{(-1_n)})\le 0$ for all $s\ge 2$ and $x\in\{0,1\}^n$. Moreover, $\sums(\interleave{x}{(-1_n)})= 0 \iff \forall i\in [n], x_i= 1$. Similarly, we have that $\sums(x)\ge 0$ and $\sums(x) = 0\iff \forall i \in [n], x_i = 0$. In other words, we have  
	\begin{itemize}
		\item $\AND(x) = \ind[\sums(\interleave{x}{(-1_n)}\ge 0 )] = \ind[\inner{\interleave{x}{1_n}}{\interleave{1_n}{(-1_n)}}_s+1 >0]$;
		\item $\OR(x) = \ind[\sums(x'_i)>0] = \ind[\inner{\interleave{x}{1_n}}{\interleave{1_n}{(0_n)}}_s >0]$.
	\end{itemize}
	Therefore for $\AND$, we can set $\theta^\AND\triangleq (W_1^\AND,W_2^\AND,b_1^\AND)$ with $W_1^{\AND}\triangleq \begin{bmatrix}
		\interleave{1_n}{(-1_n)} \\
		\interleave{1_n}{(-1_n)}
	\end{bmatrix} , b_1 = \begin{bmatrix}
		1\\
		0
	\end{bmatrix}, W_2^\AND = [1,-1]$, and we have that 
	\begin{align}
		\FF_{\theta^\AND}(\interleave{x}{1_n}) = &\rds{\relu(\inner{\interleave{x}{1_n}}{\interleave{1_n}{(-1_n)}}_s+1) - \relu(\inner{\interleave{x}{1_n}}{\interleave{1_n}{(-1_n)}}_s)} \notag\\
		= &  \ind[\inner{\interleave{x}{1_n}}{\interleave{1_n}{(-1_n)}}_s+1>0] \qquad\qquad\qquad\text{(by \Cref{lem:indicator})} \notag\\
		= & \AND(x)\notag
	\end{align}
	
	Similarly for $\OR$, we can set $\theta^\OR\triangleq (W_1^\OR,W_2^\OR,b_1^\OR)$ with $W_1^{\OR}\triangleq \begin{bmatrix}
		\interleave{1_n}{0_n} \\
		\interleave{1_n}{0_n}
	\end{bmatrix} , b_1 = \begin{bmatrix}
		0\\
		-1
	\end{bmatrix}, W_2^\OR = [1,-1]$, and we have that 
	\begin{align}
		\FF_{\theta^\OR}(\interleave{x}{1_n}) = &\rds{\relu(\inner{\interleave{x}{1_n}}{\interleave{1_n}{0_n}}_s) - \relu(\inner{\interleave{x}{1_n}}{\interleave{1_n}{0_n}}_s-1)} \notag\\
		= &  \ind[\inner{\interleave{x}{1_n}}{\interleave{1_n}{0_n}}_s>0] \qquad\qquad\qquad\text{(by \Cref{lem:indicator})} \notag\\
		= & \OR(x)\notag
	\end{align}
	The proofs for $\AND$ and $\OR$ are thus completed.
	
	Next we deal with $\MAJORITY$. Note that for $s(n)\ge \log_2 n +1$, we have that $\sum_{i=1}^n (2x_i-1) = \inner{\interleave{x}{1_n}}{\interleave{2_n}{(-1_n)}}_s $ for all $x\in\{0,1\}^n$. 
	\begin{align}
	\MAJORITY(x) = &\ind[\sum_{i=1}^n (2x_i-1)>0]  = \ind[\inner{\interleave{x}{1_n}}{\interleave{2_n}{(-1_n)}}_s>0]\notag \\
	= & \rds{\relu(\inner{\interleave{x}{1_n}}{\interleave{2_n}{(-1_n)}}_s) - \relu(\inner{\interleave{x}{1_n}}{\interleave{2_n}{(-1_n)}}_s-1)} \notag\\
	= & \FF_{\theta^\MAJORITY}(\interleave{x}{1_n}),
	\end{align}
where  $\theta^\MAJORITY\triangleq (W_1^\MAJORITY,W_2^\MAJORITY,b_1^\MAJORITY)$ with $W_1^{\MAJORITY}\triangleq \begin{bmatrix}
		\interleave{2_n}{-1_n} \\
		\interleave{2_n}{-1_n}
	\end{bmatrix} , b_1 = \begin{bmatrix}
		0\\
		-1
	\end{bmatrix}, W_2^\MAJORITY = [1,-1]$.
\end{proof}


	
%\begin{lemma}\label{lem:merge_transformer}
%For any finite vocabulary $\gV$, suppose $k\in\mathbb{N}^+$ and $\theta_1,\ldots\theta_k$	 are the parameters of $k$ decoder-only transformers with depth at most $L$, embedding width at most $d$, and precision $s$, we recursively define
%\begin{align}
%	 x_{n+i}\triangleq \transformer_{\theta_i}(x_1,\ldots, x_{n+i-1}), \quad \forall i \in [k],
%\end{align}
%There exist a decoder-only transformers with depth at most $kL+1$, embedding width at most $k(d+1)$, and precision $s+1$, whose parameter is denoted by $\theta_*$, satisfying that 
%\begin{align}
%	x_{n+k} = \transformer^k_{\theta_*}(x_1,\ldots,x_n).
%\end{align}
%\end{lemma}
%
%\begin{proof}[Proof of \Cref{lem:merge_transformer}]
%Without loss of generality, we can assume all $\{\theta_i\}_{i=1}^k$ have depths $L$ and embedding size $d$. Given the parameters $\theta_i = (\theta_{\posencoding,i},\theta_{\tokenembedding,i},\theta_{\transoutput,i},\{\theta^{(l)}_{\Attn,i},\theta^{(l)}_{\FF,i}\}_{l=0}^{L-1})$ for all $i\in[k]$. We construct the new transformer parameter $\theta_* = (\theta_{\posencoding,*},\theta_{\tokenembedding,*},\theta_{\transoutput,*},\{\theta^{(l)}_{\Attn,*},\theta^{(l)}_{\FF,*}\}_{l=0}^{kL})$ in the following way:
%\begin{itemize}
%	\item For token embedding, we set $\theta_{\tokenembedding,*}\triangleq [\theta_{\tokenembedding,1}^\top,\ldots, \theta_{\tokenembedding,k}^\top]^\top \in\mathbb{R}^{kd\times |\gV|}$;
%	\item For position encoding, we set $\theta_{\posencoding,*}\triangleq [\theta_{\posencoding,1}^\top,\ldots, \theta_{\posencoding,k}^\top]^\top \in\mathbb{R}^{kd\times n_{\max}}$, where $n_{\max}$ is the maximal sequence length satisfying $n_{\max}\ge n+k$;
%	\item For output layer, we set $\theta_{\transoutput,*}\triangleq [\theta_{\transoutput,1},\ldots, \theta_{\transoutput,k}] \in\mathbb{R}^{|\gV|\times kd}$;
%\end{itemize}
%\end{proof}




\begin{lemma}\label{lem:TC0_not_in_fixed_poly}
For all $k\in\mathbb{N}$, $\AC^0\not \subseteq \SIZE[n^k]$.	
\end{lemma}

\begin{proof}[Proof of \Cref{lem:TC0_not_in_fixed_poly}]
	We first define $\overline\SIZE[T(n)]$ as the problems solvable by circuits with  $T(n)$ standard $\AND,\OR,\NOT$ gates exactly. Thus $\SIZE[n^k] = \cup_{C\in\mathbb{N}^+} \overline{\SIZE}(Cn^k)$. Now we claim that for each $C\in\mathbb{N}$, there is a $N\in\mathbb{N}^+$, such that for all $n\ge N$, it holds that there is a conjunction normal form (CNF) with at most $n^{k+1}$ clauses over $\{x_1,\ldots,x_n\}$ that cannot be expressed by any circuit of size $Cn^k$. This claim holds because of a simple counting argument. There are at least $2^{n^{k+1}}$ different such CNFs. On the other hand, it is well known that one can represent a $T(n)$-size circuit only allowing standard $\AND,\NOT,\OR$ gates with $3T(n)\log T(n)$ bits (we need $\log T(n)$ bits to encode the id of a gate). Thus the total number of different circuits of size at most $Cn^k$ is at most $2^{3C n^k (k\log n+C)}$, which is smaller than $2^{n^{k+1}}$ for sufficiently large $n$. We denote such $n$ for each $C$ by $N_C$. Now we define the following language $\gL_{\mathsf{CNF}}$: if the input length of $x$ is $N_C$ for some $C$, use the $n^{k+1}$-clause CNF's output which cannot be expressed by size-$Cn^k$ circuits as the output; otherwise rejects (output $0$). Then clearly $\gL_{\mathsf{CNF}}\notin \overline{\SIZE}(Cn^k)$ for all $C$, thus $\gL_{\mathsf{CNF}}\notin \cup_{C\in\mathbb{N}^+} \overline{\SIZE}(Cn^k) = \SIZE[n^k]$. By construction, $\gL_{\mathsf{CNF}}\in \AC^0$. This completes the proof.
\end{proof}


%%%%% zhiyuan:commented for this verion
%\subsection{Lower Bound on }
%\zhiyuan{discussion on a specific implementation of the exponentiation function}
%\begin{lemma}\label{lem:exp_TC0}
%	For any $s(n)=O(\poly(n))$, the rounded exponential function $[\exp(\cdot)]_{s(n)}:\Floating_{s(n)}\to \Floating_{s(n)}$ can be computed by $\TC^0$ circuits, where each number in $\Floating_s$ is encoded by a binary string $a = (a_0, a_1, \ldots a_{2s})\in \{0,1\}^{2s+1}$ with value $(2a_0-1) \sum_{i=1}^{2s}2^{i-s-1}a_i$.	
%\end{lemma}
%
%\begin{proof}[Proof of \Cref{lem:exp_TC0}]
%	Given input $x =(2a_0-1) \sum_{i=1}^{2s}2^{i-s-1}a_i$, we proceed in three cases: (1) if $x\ge  s(n) $, then $\exp(x)\ge  2^{e(m)}$ and thus $[\exp(x)]_{s(n)} = 2^s- 2^{-s}$; (2) if $x\le -s(n) \cdot \ln 2$, then $\exp(x)\le  2^{-s(n)-1}$ for $s(n)\ge 2$ and thus $[\exp(x)]_{s(n)} = 0$; (3) if $|x|\le  s(n) $, by Taylor expansion of $\exp$ at $0$, we have that
%	\begin{align}
%		\abs{\exp(x) - \sum_{i=0}^{2s(n)-1} \frac{x^n}{n!}}\le \frac{(s(n))^{2s(n)}}{(2s(n))!}.
%	\end{align}
%	By Stirling's formula, we know $\sqrt{2\pi n}(\frac{n}{e})^n/n! = 1 + O(1/n)$, thus $\abs{\exp(x) - \sum_{i=0}^{2s(n)-1} \frac{x^n}{n!}}\le O(2^{-2s(n)})$
%
%\end{proof}

%%link: https://sourceware.org/git/?p=glibc.git;a=blob_plain;f=sysdeps/ieee754/ldbl-128/e_expl.c;hb=HEAD

\section{Discussion on Variants in Transformer Architecture}\label{appsec:discussion}

\subsection{Extension to Transformers with LayerNorm}\label{appsec:layernorm}

Allowing LayerNorm changes the function class that a transformer can express and the position of the layer norm also matters~\citep{xiong2020layer}. However, the expressiveness results mentioned in this work still hold for the two most popular transformer architecture variants with LayerNorm --- Post LayerNorm and Pre LayerNorm. The upper bounds on transformer expressiveness \Cref{thm:AC0_upper_bound,thm:TC0_upper_bound} clearly don't get affected by adding LayerNorm, which can be computed in polynomial time for each token. 

Below we focus on the upper bound of the expressiveness of decoder-only transformers with or without CoT. In detail, we will explain why \Cref{thm:cot_dlogn_main,thm:cot_dpolyn_slogn_main} still holds even with LayerNorm. Here the key observation is that, if each coordinate of $h\in\mathbb{R}^d$ ranges from $\{-1,1\}$ and $-1,1$ appear in pairs, then $\textsf{LayerNorm}(h) = h$. Thus it suffices to show that we can slightly twist the construction of transformers in \Cref{thm:cot_dlogn_main,thm:cot_dpolyn_slogn_main} that for all $i\in[n+T(n)], l\in\{0,0.5,1,\ldots,L\}$, $h^l_i$ is composed of $-1$ and $1$ and they appear in pairs so the sum is always $0$. Note that in the current construction, each $h^l_i$ only contains $0,-1,1$. It suffices to replace each dimension with four dimensions, in the sense $0\to (1,-1,1,-1)$, $1\to (1,1,-1,-1)$ and $-1\to (-1,-1,1,1)$. This can be done by changing the weights of the token embedding, position encoding, and the weights of the second layer of each fully-connected layer. For the outgoing layer, we just use the average of the new representation, which is exactly the same as the original value in all three cases. 

 \subsection{Extension to Transformers with Multihead Attention}\label{appsec:multihead}
 
 In this paper, for simplicity, we only focus on the case where there is only one attention head in each layer. The main results in this paper still apply if we allow constantly many attention heads, because we can simulate an attention layer with $k$ heads with $k$ attention layers with one head. Allowing an arbitrary number of attention heads while fixing total embedding size might make the constant-depth transformers strictly more expressive in certain settings and we leave it for future works.

